% Integrated working draft; autogenerated from verified standalone chapter TeX on 2026-11-12.
\section{研究背景与相关工作}
\label{sec:background}

\subsection{空间域识别的研究背景}
空间转录组技术在记录基因表达的同时保留了组织位置，使研究者能够从空间维度考察组织结构与分子异质性。与仅根据表达相似性对观测单位进行聚类不同，空间域识别（spatial domain identification）还需要考虑局部空间连续性及不同组织区域之间的差异。空间域聚类所得到的簇并不必然与组织学边界一一对应，因而算法聚类结果与生物学参考标注之间的关系需要独立评价。

人类背外侧前额叶皮层（dorsolateral prefrontal cortex, DLPFC）具有组织分层结构，是检验空间转录组区域识别方法的常用场景。Maynard等在2021年报告了来自3位独立供体的12张10x Visium切片，其中每位供体贡献两对空间邻近切片；该研究还从基因表达模式和组织层标注出发开展了多种空间结构分析\cite{maynard2021}。因此，DLPFC适合用于观察局部组织层边界，但12张切片不能理解为12个相互独立的生物学个体。

\subsection{表达表征与空间邻域信息融合}
单细胞变分推断模型（single-cell variational inference, scVI）是一类成熟的概率深度生成模型。Lopez等提出scVI用于对测序表达数据中的噪声和技术变异进行建模，并获得供下游任务使用的低维表示\cite{lopez2018}。在本研究中，scVI作为已经公开的方法用于生成spot的5维表达表征，不涉及重新设计其生成模型结构。

在空间域聚类方面，仅依赖单个spot的表达特征，可能无法充分体现其周围组织环境。Varrone等提出CellCharter，用于识别、表征和比较空间组学数据中的细胞生态位\cite{varrone2024}。CellCharter的公开API支持将中心观测单位的特征与不同层数的邻域聚合特征拼接，并支持指定所使用的邻域层数。因此，“多尺度邻域聚合”不是本研究首次提出的思想；本研究沿用其已实现功能，分析固定特征块选择对应的实证表现。

SEDR是另一类空间表示学习框架。Xu等使用掩码自监督表示学习与变分图自编码器，联合处理表达和空间邻接信息，并将其用于空间聚类等任务\cite{xu2024}。与SEDR直接学习空间融合表示不同，本研究不重新训练一套新的图自编码器，而是在既有CellCharter聚合表示中进行可控的等维特征组合比较。本文采用的后续10张DLPFC切片的参考标注表包含从SEDR项目分析数据整理的元数据字段，但其存在于SEDR资料目录并不意味着参考标签由SEDR算法预测，更不意味着属于本研究新生成的人工真值；实际标签来源须继续追溯。

\subsection{多尺度空间聚类与现有方法的关系}
多尺度空间表征已是空间转录组分析的重要研究方向。例如，Zhang等在m2ST中结合多尺度图表示，探索不同尺度中的共性和特异性信息，并采用熵机制自适应调整尺度的重要性\cite{zhang2025}。这一工作说明，研究者不仅已经采用多尺度特征，也已经探索了尺度之间的自适应融合策略。仅使用远邻特征或调整尺度权重，不能直接构成区别于已有方法的原创算法贡献。

本文所进行的等维特征块实验，重点不在于提出新一代尺度融合网络，而在于控制输入维度与聚类设置后，观察相同CellCharter表示中不同邻域层组合对应的性能变化。具体而言，将已生成的25维多尺度向量划分为5个各5维的特征块，比较对照组合$L_0+L_1+L_2$与候选组合$L_0+L_1+L_4$。两组均为15维，分别使用相同配置的GMM聚类。在没有跨切片和跨供体稳定优势的证据之前，这一比较应表述为受控消融与实证分析，而非经过验证的新方法。

\subsection{空间聚类评估、边界指标与参考标注局限}
常规空间域聚类评价通常采用调整兰德指数（adjusted Rand index, ARI）、归一化互信息（normalized mutual information, NMI）等指标衡量预测簇与参考层标签的一致性。但是，整体聚类指标不能直接解释模型在哪类组织交界处检出了边界，也不能区分错误边界减少与真实边界漏检增加所带来的影响。

近年来针对空间聚类基准的研究还指出，性能报告高度依赖数据处理、模型配置及参考标注。Sun等于2026年在\textit{Nature Methods}中提出SACCELERATOR框架，系统讨论传统空间聚类基准中模型表现差异、参考标注不确定性以及专家参与评价的必要性\cite{sun2026}。这一发现尤其提醒研究者，DLPFC组织层标签应被视为可追溯但可能存在误差和粒度差异的参考，而不能被无条件当作绝对真值。

基于上述认识，本文除报告ARI、NMI外，还在Visium六邻接图上定义空间边界F1，并分别统计正确检出的层间边界（TP）、错误预测的边界（FP）及漏检的参考边界（FN）。对于选定的正向与负向说明性案例，进一步按参考组织层对计算真实边界的召回率。两张案例是在已观察整体指标后选取，仅用于解释不同边界的检出变化，不作为独立随机测试。该评价是本研究实际实施的分析步骤；在未完成针对全部相关评价文献的系统性近邻核查前，不宣称首次提出边界F1或按层对的召回率。

\subsection{研究问题与工作定位}
综上，本研究围绕以下三个可以直接由已有实验回答的问题开展分析：
\begin{enumerate}
  \item 当输入维度与GMM配置保持一致时，改变CellCharter邻域特征块是否能在DLPFC切片上获得一致的聚类性能改善？
  \item 当总体边界F1发生变化时，真实边界检出数量与错误预测边界数量分别发生了怎样的变化？
  \item 在典型切片中，不同参考组织层之间的边界召回率是否呈现一致变化？
\end{enumerate}
本研究的主要工作是基于成熟CellCharter/scVI框架建立可复现的等维特征对照、完成12张切片与3位供体的描述性比较，并进一步分析部分组织层间边界的误报与漏检权衡。研究的结论以实际实验数值为依据，不预设候选特征组合具有普适优势，也不将已有多尺度方法的组成部分重新命名为新算法。
